\subsection{Eigene Rolle im Unternehmen und im Projekt}
\label{subsec:EigeneRolle}

\par Als \gls{pm} innerhalb des Bereichs \gls{itt} wurde mir die Position des \gls{wsl} für den Bereich \gls{ricefw} übertragen. In dieser Schlüsselfunktion für \gls{gls-erp}-Systeme verantwortete ich die IT-seitige Umsetzung sämtlicher kundenindividueller Entwicklungsobjekte, Schnittstellen und komplexen Workflows. Meine Rolle zeichnet sich durch eine hohe Management-Verantwortung in einer komplexen Matrixorganisation aus: Ich führte lateral ein Team aus internen Mitarbeitenden der Finanz- und IT-Abteilungen sowie externen Experten aus sieben verschiedenen Zulieferunternehmen. Über die fachliche Steuerung hinaus oblag mir die Budgetverantwortung innerhalb des Workstreams sowie das Reporting an die technische Programmleitung und den internen Führungskreis der \gls{ais}. Die Position erforderte ein hohes Maß an Stakeholder-Management, um die Interessen der verschiedenen Unternehmenseinheiten mit den technischen Lieferzielen zu synchronisieren.



\begin{comment}
\par In der Stammorganisation der \ac{ais} bin ich dem Bereich \acl{itt} (\ac{itt}) zugeordnet, in dem ich als \acl{pm} (\ac{pm}) tätig bin. Der Bereich stellt die zentrale Abteilung für die IT-\ac{pm} der \ac{ais} dar. Die IT-\acl{pl} (\ac{pl}) und erfolgskritischen Positionen der laufenden IT-Groß- und Transformationsprogramme \ac{r2h} sowie \acl{vita} (\acs{vita}) werden durch Mitarbeitende des \ac{itt} besetzt. Hinzu kommen weitere IT-Projekte wie die Portierung der ADAC-Webseite von Version 3.0 auf 4.0, die ebenfalls durch \ac{itt} begleitet bzw. geleitet werden.

\par Seit meinem Eintritt in \ac{itt} umfasst mein primäres Einsatzgebiet das Programm \ac{r2h} sowie eine beratende Funktion zur Umgebungs- und Releaseplanung im Projekt \acl{inex}. Im Programm \ac{r2h} bin ich als leitender \acs{pm} bzw. \acl{wsl} (\ac{wsl}) für die IT-Umsetzung aller \acl{ricefw}-Objekte (\acs{ricefw}) verantwortlich.\footnote{\href{https://de.wikipedia.org/wiki/RICEFW}{\ac{ricefw}} bezeichnet im SAP-Kontext die sechs zentralen Entwicklungsobjekte – Reports/ Berichte also individuelle Auswertungen und Datendarstellungen - Interfaces / Schnittstellen als Drittsystemanbindungen - Conversions/ Konvertierungen also Programme zur Datenmigration und Datenformatkonvertierungen - Enhancements/ Erweiterungen von Standardfunktionalität - Forms /Formulare also Druckausgaben und Layouts - Workflows/ automatisierte Geschäftsprozesse.} Neben der Implementierung im \ac{s4} entsteht Anpassungsbedarf im \ac{r3}, um den temporären parallelen Betrieb beider Systeme, welche sich durch die nach \acl{bk} (\ac{bk}) gestückelten Migrationstrategie bedingt abzubilden.\footnote{Ein \href{https://de.wikipedia.org/wiki/Rechnungskreis}{\acl{bk}} (\ac{bk}) ist die kleinste organisatorische Einheit des externen Rechnungswesens im SAP.}

\par Zu meinen Aufgaben gehört die Planung und Steuerung der Liefergegenstände und des Budgets, die Überwachung des Projektfortschritts, die Abstimmung mit sieben fachlichen \ac{wsl} sowie die laterale Führung zahlreicher Umsetzender aus unterschiedlichen Zulieferunternehmen zusätzlich zum Hauptimplementierungspartner. Darüber hinaus fungiere ich als Stellvertretung der technischen Programmleitung und berichte sowohl an die Gesamtprogrammleitung als auch in das wöchentliche Transformations-Statusmeeting sowie den internen Führungskreis der \ac{ais}.

\par Der \ac{ws} \ac{ricefw} ist in einer Matrixorganisation mit internen Mitarbeitenden der Finanzabteilung (\ac{fnz}) des \ac{adac}, der \ac{ais} sowie externen Zulieferern aus insgesamt sieben Unternehmen besetzt. Die Mitarbeitenden sind entweder vollständig, stundenweise oder ad-hoc für das Projekt verfügbar. Lateral bin ich ihnen weisungsbefugt.
\end{comment}